\chapter*{แผนบริหารการสอนประจำวิชา}
\addcontentsline{toc}{chapter}{แผนบริหารการสอนประจำวิชา}

\section*{ข้อมูลทั่วไปของรายวิชา}
\begin{itemize}
    \item \textbf{รหัสและชื่อรายวิชา:} ฟิสิกส์ 2026 คณิตศาสตร์สำหรับฟิสิกส์ (Mathematics for Physics)
    \item \textbf{จำนวนหน่วยกิต:} 3 หน่วยกิต (3-0-6)
    \item \textbf{หลักสูตรและประเภทของรายวิชา:} วิทยาศาสตรบัณฑิต สาขาวิชาฟิสิกส์ หมวดวิชาเฉพาะด้าน
    \item \textbf{อาจารย์ผู้รับผิดชอบรายวิชาและผู้สอน:} ดร.ชีวะ ทัศนา
    \item \textbf{ภาคการศึกษา / ชั้นปีที่เรียน:} ภาคการศึกษาที่ 1 ชั้นปีที่ 2
\end{itemize}

\section*{คำอธิบายรายวิชา (Course Description)}
ศึกษาและฝึกทักษะการคำนวณทางคณิตศาสตร์ขั้นสูงสำหรับประยุกต์ใช้ในฟิสิกส์ ได้แก่ ระบบพิกัดคาร์ทีเซียน พิกัดทรงกระบอก และพิกัดทรงกลม การวิเคราะห์เวกเตอร์ เกรเดียนต์ ไดเวอร์เจนซ์ เคิร์ล อนุกรมกำลัง อนุกรมเทย์เลอร์ การประมาณค่าในทางฟิสิกส์ จำนวนเชิงซ้อนและการสั่นฮาร์มอนิก เมทริกซ์และพีชคณิตเชิงเส้น ค่าเจาะจงและเวกเตอร์เจาะจง แคลคูลัสเชิงอนุพันธ์หลายตัวแปร ตัวคูณลากรองจ์ ปริพันธ์หลายชั้น ทฤษฎีบทของเกาส์ ทฤษฎีบทของสโตกส์ ตลอดจนสมการเชิงอนุพันธ์สามัญอันดับหนึ่งและอันดับสองในการจำลองระบบกายภาพ

\section*{วัตถุประสงค์การเรียนรู้ (Course Learning Outcomes: CLOs)}
เมื่อสิ้นสุดการเรียนการสอนรายวิชานี้ นักศึกษาสามารถ:
\begin{enumerate}
    \item \textbf{CLO 1:} อธิบายมโนทัศน์และหลักการทางคณิตศาสตร์ที่ใช้ในกลศาสตร์ แม่เหล็กไฟฟ้า และควอนตัมฟิสิกส์ได้อย่างถูกต้อง
    \item \textbf{CLO 2:} เลือกใช้ระบบพิกัดและตัวดำเนินการเชิงอนุพันธ์เวกเตอร์ในการตั้งสมการและแก้ปัญหาฟิสิกส์ได้อย่างเหมาะสม
    \item \textbf{CLO 3:} แก้สมการเชิงอนุพันธ์สามัญอันดับหนึ่งและสอง เพื่อวิเคราะห์พฤติกรรมการเคลื่อนที่และการสั่นของระบบกายภาพได้
    \item \textbf{CLO 4:} เชื่อมโยงระเบียบวิธีทางคณิตศาสตร์เข้ากับการจำลองสถานการณ์ทางฟิสิกส์จริงได้อย่างมีวิจารณญาณ
\end{enumerate}

\section*{การวัดและประเมินผลการเรียนรู้}
\begin{itemize}
    \item การประเมินระหว่างภาคเรียน (แบบฝึกหัดท้ายบทและโครงงานย่อย): ร้อยละ 60
    \item การสอบวัดผลปลายภาคเรียน: ร้อยละ 40
\end{itemize}
